% Part: proof-theory
% Chapter: proof-search
% Section: introduction
%READERNOTE{SOL6-A314: पुढील पद-प्रतिमानावर आधारित संपूर्णतेची चर्चा वाक्यांच्या क्रमवर्तींसाठी आहे, ही व्याप्ती परिचयात स्पष्ट केली आहे.}

\documentclass[../../../include/open-logic-section]{subfiles}

\begin{document}

\olfileid{pt}{ps}{int}

\olsection{प्रस्तावना}

\Log{G3c} किंवा~\Log{N1i} सारख्या !!{proof}
पद्धतींची स्वयंसिद्धके आणि अनुमाननियम ठरवतात की
क्रमवर्ती किंवा !!{formula}s यांच्या कोणत्या वृक्षांना
योग्य !!{proof}s म्हणायचे. असा क्रमवर्तींचा किंवा
!!{formula}s चा वृक्ष दिला असेल (आवश्यक असल्यास
प्रत्येक टप्प्यावर कोणता नियम लागू झाला, किंवा
गृहीतकांच्या आणि त्यांना निरस्त करणाऱ्या अनुमाननियमांच्या
खुणा यांसारखी अतिरिक्त माहितीही असेल), तर
स्वयंसिद्धके आणि नियम यांच्या व्याख्यांवरून तो
वृक्ष योग्य !!{proof} आहे का हे तपासता येते.
पण दिलेल्या क्रमवर्तीची किंवा दिलेल्या गृहीतकांपासून
एखाद्या !!{formula} ची !!a{proof} \emph{शोधणे}
हा वेगळा आणि बहुतेकदा अधिक कठीण प्रश्न आहे.
हाच \emph{सिद्धता-शोध} प्रश्न.

सिद्धता-शोधासाठी एक अगदी सोपी पद्धत आहे.
!!{formula}s यांना सांत वर्णमालेतील चिन्हांचे अनुक्रम
समजू शकतो. क्रमवर्ती आणि क्रमवर्तींचे किंवा
!!{formula}s चे वृक्षही !!{formula}s चे अनुक्रम
मानता येतात (वृक्षाची शाखा दाखवण्यासाठी विशेष
कंस वापरून). चिन्हांचा एखादा अनुक्रम योग्य !!{proof}
दर्शवतो का, हे स्वयंसिद्धके आणि नियम यांमुळे
यांत्रिकरीत्या तपासता येते. परिणामी, सर्व \emph{शक्य}
योग्य !!{proof}s यांत्रिकरीत्या निर्माण करता येतात:
!!{proof}s दर्शवणाऱ्या वर्णमालेतील चिन्हांचे सर्व
अनुक्रम तयार करून प्रत्येक उमेदवार अनुक्रम योग्य
सिद्धता आहे का ते तपासा. ही सोपी सिद्धता-शोध
पद्धत अशी: दिलेल्या क्रमवर्तीची !!{proof} मिळेपर्यंत
सर्व योग्य !!{proof}s निर्माण करा. (किंवा दिलेल्या
गृहीतकांपैकी काही मुक्त गृहीतके असलेल्या दिलेल्या
!!{formula} ची !!{proof} मिळेपर्यंत तसे करा.)

याहून चांगली पद्धत म्हणजे हाताने !!a{proof}
शोधताना वापरायची सहज पद्धत: अंतिम क्रमवर्तीपासून
सुरुवात करा, एक !!a{formula} निवडा, आणि अनुमाननियम
``उलट्या'' दिशेने लावून एक किंवा अनेक आधारविधाने
निर्माण करा. त्यांची !!{proof}s असतील, तर शेवटच्या
अनुमानासह त्यांतून दिलेल्या क्रमवर्तीची !!a{proof}
बनते. नैसर्गिक निगमन पद्धतीत !!a{proof}
शोधण्यासाठीही अशीच, पण अधिक गुंतागुंतीची, पद्धत
वापरली जाईल.

मात्र ही पद्धत हमखास यशस्वी नाही:
चुकीचा नियम किंवा !!{formula}s चा चुकीचा क्रम
निवडल्यास मार्ग बंद होऊ शकतो. ती पूर्णपणे
निर्दिष्टही नाही: प्रत्येक टप्प्यावर निवडण्यासाठी
एकापेक्षा अधिक !!{formula} किंवा उलट्या दिशेने
लावण्यासाठी एकापेक्षा अधिक नियम असू शकतात.
\emph{सिद्धता-शोध अल्गोरिदम} म्हणजे पूर्णपणे
निर्दिष्ट अशी प्रक्रिया की !!a{proof} अस्तित्वात
असेल तर ती शोधून काढते (म्हणजे ती हमखास यशस्वी आहे).

काही !!{proof} पद्धतींसाठी सिद्धता-शोध
अल्गोरिदम इतरांपेक्षा सोपे असतात. क्रमवर्ती-पद्धतीत
!!a{formula} चा !!{main operator} आणि ते डावीकडे
की उजवीकडे आहे यावर उलट्या दिशेने कोणता नियम
लावायचा हे ठरते. उदाहरणार्थ,~$!A \land !B$
उजवीकडे असलेल्या~$\Gamma \Sequent \Delta, !A \land !B$
ची !!a{proof} शोधायची असेल, तर~\RightR{\land}
उलटा लावून अनुरूप दोन आधारविधाने
$\Gamma \Sequent \Delta, !A$ आणि
$\Gamma \Sequent \Delta, !B$ निर्माण करावीत.
पण पद्धत~\Log{G1c} असेल, तर आकुंचनामुळे अधिक
पर्याय मिळतात आणि काम कठीण होते: उलट्या दिशेने
\RightR{\Contraction} लावून
$\Gamma \Sequent \Delta, !A \land !B, !A \land !B$
हे आधारविधानही तयार करता येईल. मग
$\Gamma \Sequent \Delta, !A \land !B, !A \land !B, !A \land !B$
हे, आणि असे पुढे. \Log{G2c} मध्ये अडचण अधिक आहे.
\RightR{\land} उलटा लावताना~$\Gamma_1$, $\Gamma_2$
हे~$\Gamma = \Gamma_1 \cup \Gamma_2$ पूर्ण करतील
असा अंदाज, आणि~$\Delta_1$, $\Delta_2$ हे
$\Delta = \Delta_1 \cup \Delta_2$ पूर्ण करतील
असा अंदाज बांधून,~$\Gamma_1 \Sequent \Delta_1, !A$
आणि~$\Gamma_2 \Sequent \Delta_2, !B$ ही
आधारविधाने वापरून पाहावी लागतात. चुकीच्या अंदाजाने
पुढे मार्ग बंद झाला, तर सुरुवातीला परत जाऊन
वेगळे~$\Gamma_1$, $\Gamma_2$, $\Delta_1$, $\Delta_2$
निवडावे लागतात. !!{formula} हे~$!A \lor !B$
असेल, तर~\RightR{\lor} चे दोन पर्यायांपैकी एक
निवडून~$\Gamma \Sequent \Delta, !A$ किंवा
$\Gamma \Sequent \Delta, !B$ या दोनपैकी
एक आधारविधान निर्माण करावे लागते. चुकीच्या
निवडीवर पुन्हा मागे फिरावे लागते.

\Log{G3c} मध्ये या अडचणी नाहीत.
तिथे संरचनात्मक नियम नाहीत; !!{main operator}
आणि !!{formula} ची बाजू यांवर अनुमाननियम ठरतो.
म्हणून~\Log{G3c} हे~\Log{G1c} किंवा~\Log{G2c}
पेक्षा सिद्धता-शोधासाठी अधिक सोयीचे आहे.
यशस्वी शोधातून~\Log{G3c} मधील !!a{proof} मिळते;
नंतर त्या !!{proof} चे~\Log{G1c} किंवा~\Log{G2c}
(किंवा~\Log{N1c}) मध्ये रूपांतर करता येते.
त्यामुळे~\Log{G3c} साठी सिद्धता-शोध अल्गोरिदम
आणि~\Log{G3c} !!{proof}s चे इतर पद्धतींतील
!!{proof}s मध्ये रूपांतर करणारे अल्गोरिदम मिळाले,
तर त्या इतर पद्धतींचा सिद्धता-शोध प्रश्नही सुटतो.

सिद्धता-शोध अल्गोरिदम हमखास यशस्वी आहे
हे कसे समजेल? त्याने !!a{proof} शोधली नाही,
तर !!{proof} अस्तित्वातच नाही हे दाखवावे लागेल.
कारण अल्गोरिदम चुकीचा निवडला, तर सिद्धता
असूनही ती न सापडण्याचे प्रसंग येऊ शकतात.
अगदी सोप्या पद्धतीत ही अडचण नसते: ती सर्व
शक्य !!{proof}s तपासते. पण सिद्धता-शोध
अल्गोरिदम कार्यक्षम असावा आणि शक्य तितके कमी
प्रयत्न वापरावेत अशी अपेक्षा असते.

पुढील बंद-पद शोध आणि त्याच्या पद-प्रतिमानाची सिद्धता
वाक्यांच्या क्रमवर्तींसाठी मांडली आहे; खालील पूर्तिचे
वर्णन या व्याप्तीत वाचायचे.
सिद्धता-शोध अल्गोरिदम हमखास यशस्वी आहे
हे दाखवण्याची मुख्य कल्पना अशी: शोध अयशस्वी
झाल्यास त्याच्या अपयशाच्या स्वरूपावरून त्या
क्रमवर्तीची !!a{proof} अस्तित्वात असू शकत नाही हे दाखवा.
हे दाखवण्यासाठी ती वैध नाही असे दाखवू.
अभिजात प्रथम-क्रम तर्कशास्त्रात
$\Gamma \Sequent \Delta$ ही क्रमवर्ती वैध
असण्याचा अर्थ असा की प्रत्येक !!{structure}
मध्ये~$\Gamma$ मधील किमान एक !!{formula}
असत्य असेल किंवा~$\Delta$ मधील किमान एक
!!{formula} सत्य असेल. \Log{G3c} \emph{निर्दोष}
आहे; म्हणजे तो फक्त वैध क्रमवर्ती !!{prove}s करतो.
हे !!{proof}s वर विगमन करून सिद्ध होते:
स्वयंसिद्धके स्पष्टपणे वैध आहेत आणि नियमांची
आधारविधाने वैध असतील तर निष्कर्षही वैध असतो.
सिद्धता-शोध अल्गोरिदम हमखास यशस्वी आहे हे
दाखवण्यासाठी,~$\Gamma \Sequent \Delta$ ची
!!a{proof} शोधण्यात अपयश आल्यास अशी
!!a{structure} परिभाषित करता येते हे दाखवतो
की तिच्यात~$\Gamma$ मधील सर्व !!{formula}s
सत्य आणि~$\Delta$ मधील सर्व !!{formula}s
असत्य आहेत. यावरून~\Log{G3c} \emph{संपूर्ण}
आहे हेही सिद्ध होते:~$\Gamma \Sequent \Delta$
वैध असेल, तर तिची~\Log{G3c}-!!{proof} असते.
अयशस्वी सिद्धता-शोधामुळे~$\Gamma \Sequent \Delta$
अवैध ठरते; म्हणून वैध क्रमवर्तीसाठीचा प्रत्येक
सिद्धता-शोध यशस्वी ठरलाच पाहिजे.

\end{document}
